\section*{Chapitre \ref{ch:g6:solutions-and-solubility} --- Solutions et solubilité}

\begin{solution}{exo:g6:solutions-and-solubility:1}
\omterm{def:g6:solutions-and-solubility:solute}{Soluté}~: le sucre. \omterm{def:g6:solutions-and-solubility:solute}{Solvant}~: l'eau. Oui, le \omterm{def:g6:solutions-and-solubility:solute}{solvant} est l'eau~: c'est
une \omterm{def:g6:solutions-and-solubility:solute}{solution aqueuse}.
\end{solution}

\begin{solution}{exo:g6:solutions-and-solubility:2}
Une solution qui ne peut plus dissoudre davantage de son \omterm{def:g6:solutions-and-solubility:solute}{soluté}. Le
\omterm{def:g6:solutions-and-solubility:solute}{soluté} qu'on y ajoute reste au fond sans \omterm{def:g2:mixing-and-dissolving:dissolve}{se dissoudre}.
\end{solution}

\begin{solution}{exo:g6:solutions-and-solubility:3}
$250 + 20 = \qty{270}{g}$.
\end{solution}

\begin{solution}{exo:g6:solutions-and-solubility:4}
L'éthanol~: \omterm{def:g6:solutions-and-solubility:miscible}{miscible}. L'huile de cuisine~: \omterm{def:g6:solutions-and-solubility:miscible}{non miscible}. Le vinaigre~:
\omterm{def:g6:solutions-and-solubility:miscible}{miscible} (c'est surtout de l'eau).
\end{solution}

\begin{solution}{exo:g6:solutions-and-solubility:5}
La propanone (acétone). GHS02~: très inflammable~; GHS07~: irrite les
yeux, la vapeur peut provoquer somnolence ou vertiges.
\end{solution}

\begin{solution}{exo:g6:solutions-and-solubility:6}
$360 \times \frac{500}{1000} = \qty{180}{g}$.
\end{solution}

\begin{solution}{exo:g6:solutions-and-solubility:7}
\qty{200}{g} d'eau dissolvent au plus $360 \times \frac{200}{1000} =
\qty{72}{g}$. Non~: $80 - 72 = \qty{8}{g}$ restent au fond.
\end{solution}

\begin{solution}{exo:g6:solutions-and-solubility:8}
La solution pèse $225 + 25 = \qty{250}{g}$~; $\frac{25}{250} =
\qty{10}{\percent}$.
\end{solution}

\begin{solution}{exo:g6:solutions-and-solubility:9}
Dans le premier bécher, où tout le sel s'est dissous et où la solution
n'est pas encore saturée.
\end{solution}

\begin{solution}{exo:g6:solutions-and-solubility:10}
Le dioxyde de carbone a été dissous sous pression dans la bouteille
fermée. L'ouvrir fait baisser la pression~: l'eau peut contenir moins de
gaz, et le gaz en trop ressort sous forme de bulles.
\end{solution}

\begin{solution}{exo:g6:solutions-and-solubility:11}
$2000 \div 0.04 = 50\,000$ fois plus.
\end{solution}

\begin{solution}{exo:g6:solutions-and-solubility:12}
Un gaz \omterm{def:g2:mixing-and-dissolving:dissolve}{se dissout} moins dans l'eau chaude~: l'étang réchauffé contient
moins d'oxygène dissous, et les poissons manquent d'oxygène.
\end{solution}

\begin{solution}{pb:g6:solutions-and-solubility:1}
\textbf{1.} \omterm{def:g6:solutions-and-solubility:solute}{Soluté}~: le sel. \omterm{def:g6:solutions-and-solubility:solute}{Solvant}~: l'eau.

\textbf{2.} L'eau contient autant de sel qu'elle le peut~: plus aucun sel
ne peut s'y dissoudre.

\textbf{3.} Le sel non dissous au fond prouve que la saumure est
saturée~: elle contient le plus de sel possible.

\textbf{4.} $360 \times 20 = \qty{7200}{g} = \qty{7.2}{kg}$.

\textbf{5.} $20 + 7.2 = \qty{27.2}{kg}$.

\textbf{6.} $\frac{7.2}{27.2} \approx 0.26$~: environ \qty{26}{\percent}.

\textbf{7.} Non~: l'eau pourrait dissoudre \qty{7.2}{kg}. Elle pourrait
encore en dissoudre $7.2 - 5 = \qty{2.2}{kg}$.

\textbf{8.} $20 - 5 = \qty{15}{kg}$ d'eau (\qty{15}{L}).

\textbf{9.} $360 \times 15 = \qty{5400}{g} = \qty{5.4}{kg}$.

\textbf{10.} Il sort de la solution sous forme de cristaux solides~: il
cristallise et tombe au fond.

\textbf{11.} De nouveaux cristaux blancs apparaissent et s'entassent au
fond de la cuve.

\textbf{12.} $7.2 - 5.4 = \qty{1.8}{kg}$ de sel ont cristallisé.
\end{solution}
